\documentclass[11pt,a4paper]{article}

\usepackage[T1]{fontenc}
\usepackage[utf8]{inputenc}
\IfFileExists{lmodern.sty}{\usepackage{lmodern}}{}
\IfFileExists{microtype.sty}{\usepackage{microtype}}{}
\usepackage{amsmath,amssymb,mathtools}
\usepackage{booktabs}
\usepackage{array}
\usepackage{enumitem}
\usepackage[margin=25mm]{geometry}
\usepackage{xcolor}
\usepackage{listings}
\usepackage[hidelinks]{hyperref}

\definecolor{codebg}{HTML}{F4F6F8}
\definecolor{codekey}{HTML}{1F4E79}
\definecolor{codecomment}{HTML}{607D8B}
\lstdefinelanguage{LayoutJSON}{
  basicstyle=\ttfamily\small,
  backgroundcolor=\color{codebg},
  frame=single,
  rulecolor=\color{black!15},
  breaklines=true,
  showstringspaces=false,
  columns=fullflexible,
  string=[b]",
  stringstyle=\color{codekey},
  comment=[l]{//},
  commentstyle=\color{codecomment},
}
\lstset{language=LayoutJSON}

\setlist[itemize]{leftmargin=*,itemsep=0.25em}
\setlist[enumerate]{leftmargin=*,itemsep=0.25em}
\setlength{\parindent}{0pt}
\setlength{\parskip}{0.55em}

\newcommand{\R}{\mathbb{R}}
\newcommand{\SE}{\mathrm{SE}(3)}
\newcommand{\Frame}{\mathsf{F}}
\newcommand{\Rot}{\operatorname{Rot}}
\newcommand{\Apply}{\operatorname{Apply}}
\newcommand{\Path}{\operatorname{Path}}
\newcommand{\ops}{\mathcal{O}}
\newcommand{\eps}{\varepsilon}
\newcommand{\trans}{\mathsf{T}}
\newcommand{\identity}{\mathsf{I}}

\title{A Curve-Referenced Positioning Model for Beam-Line Layouts}
\author{}
\date{Prototype specification --- 5 September 2026}

\begin{document}
\maketitle

\begin{abstract}
This document specifies a three-dimensional positioning model intended to serve
both installation geometry and beam-line modelling.  A layout contains named
reference curves, reusable types, and positioned object instances.  Every
position of interest is a right-handed frame, not merely a point.  Curves and
type geometries may be curved; object placement can align any named target frame
of an object with a transformed curve frame or with a frame on another object.
The document defines the canonical data model and the mathematics required to
obtain world frames, optional swept mechanical geometry, independent magnetic
and beam-interface axes and boundary planes, and a curve station inferred from
an arbitrary world point.
\end{abstract}

\tableofcontents

\section{Scope and design principles}

The model separates three concerns:

\begin{itemize}
  \item A \emph{reference curve} supplies a stable spatial path and a frame at
        every path coordinate.
  \item A \emph{type} owns optional reusable mechanical geometry, optional
        magnetic and beam-interface axes, and named local frames.
  \item An \emph{object} is an instance of a type.  It owns only its type name
        and a positioning rule.
\end{itemize}

This separation permits many physical instances to share one type while being
placed independently.  It also keeps the installation frame graph distinct
from a subsequently derived accelerator model.  This specification defines
geometry and placement; it does not prescribe how a MAD-X sequence or its
misalignments are fitted to that geometry.

The JSON format described here is the canonical prototype format.  There is no
legacy migration layer.  All stored lengths are in metres, curvatures in
metres$^{-1}$, and angles in radians.  A user interface may display degrees or
different step sizes without changing the stored units.

\section{Frames, coordinates, and composition}

\subsection{World and local axes}

The fixed world axes are denoted $X,Y,Z$.  A local frame is
\begin{equation}
  \Frame=(\mathbf{o},\mathbf{x},\mathbf{y},\mathbf{t}),
\end{equation}
where $\mathbf{o}\in\R^3$ is its origin and
$\mathbf{x},\mathbf{y},\mathbf{t}$ are an orthonormal, right-handed triad:
\begin{equation}
  \mathbf{x}\times\mathbf{y}=\mathbf{t}.
\end{equation}
The vector $\mathbf{t}$ is the local longitudinal axis and, on a path, its
tangent.  The data model and user interface call this axis \texttt{s}; this
document uses $s$ for the scalar path coordinate and $\mathbf{t}$ for the vector
to avoid ambiguity.

The homogeneous matrix of the frame is
\begin{equation}
  [\Frame]
  =
  \begin{bmatrix}
    \mathbf{x} & \mathbf{y} & \mathbf{t} & \mathbf{o}\\
    0&0&0&1
  \end{bmatrix}
  \in\SE.
  \label{eq:frame-matrix}
\end{equation}
It maps a local point $\mathbf{q}=(q_x,q_y,q_s)^\mathsf{T}$ to world
coordinates:
\begin{equation}
  \mathbf{p}_{W}
  =\mathbf{o}+q_x\mathbf{x}+q_y\mathbf{y}+q_s\mathbf{t}.
  \label{eq:local-to-world}
\end{equation}
The world frame is
$\identity=(\mathbf{0},\mathbf{e}_X,\mathbf{e}_Y,\mathbf{e}_Z)$.

Frame composition is ordinary matrix multiplication.  If $\Frame_A$ is a
frame in world coordinates and $\Frame_B$ is expressed locally in $A$, then
\begin{equation}
  [\Frame_{W\leftarrow B}]
  = [\Frame_{W\leftarrow A}][\Frame_{A\leftarrow B}].
  \label{eq:composition}
\end{equation}

\subsection{Ordered elementary operations}

A transformation is stored as an ordered list of two-element arrays
\texttt{[name,value]}.  Except for \texttt{ts}, operations act on the
\emph{current} local axes.  Equivalently, their elementary matrices are
post-multiplied onto the current frame.

Let
\begin{equation}
  T_x(a)=
  \begin{bmatrix}1&0&0&a\\0&1&0&0\\0&0&1&0\\0&0&0&1\end{bmatrix},\quad
  T_y(a)=
  \begin{bmatrix}1&0&0&0\\0&1&0&a\\0&0&1&0\\0&0&0&1\end{bmatrix},
\end{equation}
and
\begin{equation}
  T_t(a)=
  \begin{bmatrix}1&0&0&0\\0&1&0&0\\0&0&1&a\\0&0&0&1\end{bmatrix}.
\end{equation}
The local-axis rotation matrices are
\begin{align}
  R_x(\alpha)&=
  \begin{bmatrix}
    1&0&0\\0&\cos\alpha&-\sin\alpha\\0&\sin\alpha&\cos\alpha
  \end{bmatrix},\\
  R_y(\beta)&=
  \begin{bmatrix}
    \cos\beta&0&\sin\beta\\0&1&0\\-\sin\beta&0&\cos\beta
  \end{bmatrix},\\
  R_t(\gamma)&=
  \begin{bmatrix}
    \cos\gamma&-\sin\gamma&0\\\sin\gamma&\cos\gamma&0\\0&0&1
  \end{bmatrix}.
\end{align}
They are embedded in homogeneous matrices with zero translation.

The operations have the following meanings.

\begin{center}
\begin{tabular}{>{\ttfamily}l l l}
\toprule
\normalfont Operation & Meaning & Unit \\
\midrule
tx & translate along current $\mathbf{x}$ & m \\
ty & translate along current $\mathbf{y}$ & m \\
tt & translate in a straight line along current $\mathbf{t}$ & m \\
rx & right-handed rotation about current $\mathbf{x}$ & rad \\
ry & right-handed rotation about current $\mathbf{y}$ & rad \\
rs & right-handed rotation about current $\mathbf{t}$ & rad \\
ts & advance along a curve or the type-local mechanical path & m \\
\bottomrule
\end{tabular}
\end{center}

Thus, for a list containing no \texttt{ts},
\begin{equation}
 [\Apply(\Frame,[o_1,\ldots,o_n])]
 = [\Frame]E(o_1)\cdots E(o_n).
 \label{eq:ordered-ops}
\end{equation}
Order matters.  For example, a translation followed by a rotation generally
does not equal the same rotation followed by that translation.

The operation \texttt{tt} is an ordinary straight translation.  The operation
\texttt{ts} is \emph{not} an elementary rigid transform independent of context:
it requests movement along a specified curved path and changes the tangent
frame accordingly.  Its two evaluation modes are defined in
Sections~\ref{sec:reference-transform} and~\ref{sec:type-local}.

\section{Canonical entity model}

\subsection{Layout}

A layout is a strict JSON object with exactly three name-indexed dictionaries:
\begin{equation}
  \mathcal{L}=(\texttt{reference\_curves},\texttt{types},\texttt{objects}).
\end{equation}
Names are non-empty strings and are unique within their dictionary.  Every
symbolic reference must resolve.

\subsection{References and transformations}

A reference is exactly one of
\begin{align*}
 &\texttt{\{"kind":"world"\}},\\
 &\texttt{\{"kind":"curve","curve":curveName\}},\\
 &\texttt{\{"kind":"object\_frame","object":objectName,
             "frame":frameName\}}.
\end{align*}
The third alternative names a complete frame belonging to an object.

A general transformation has the form
\begin{lstlisting}
{
  "reference": <reference>,
  "transformation": [[<operation>, <finite number>], ...]
}
\end{lstlisting}
and is used for a curve's starting frame.  In a general transformation,
\texttt{ts} is valid only with a curve reference.

An object position extends it with a required target:
\begin{lstlisting}
{
  "target": "center or named frame",
  "reference": <reference>,
  "reference_curve": "optional curve name",
  "transformation": [[<operation>, <finite number>], ...]
}
\end{lstlisting}
The field \texttt{target} is required in the canonical JSON.  The conventional
target is \texttt{center}.  The optional \texttt{reference\_curve} has meaning
only when a world or object-frame reference is combined with \texttt{ts}; in
that case it is required and supports the station-inference rule of
Section~\ref{sec:s-inference}.  It is forbidden when the primary reference is
already a curve.  If it is present without \texttt{ts}, it has no geometric
effect and should be omitted.

\subsection{Reference curves}

A reference curve contains
\begin{lstlisting}
{
  "color": "#RRGGBB",
  "starting_frame": <transformation>,
  "segments": [[length, angle, roll], ...]
}
\end{lstlisting}
There is at least one segment.  Each segment length is strictly positive;
\texttt{angle} is the signed total bend angle and \texttt{roll} selects the bend
plane.  The color is visual metadata and does not affect geometry.

\subsection{Types}

A type is a reusable definition:
\begin{lstlisting}
{
  "shape": <optional shape tuple>,
  "color": "#RRGGBB",
  "magnetic_center": <optional local transformation>,
  "magnetic_length": <optional positive number>,
  "magnetic_curvature": <optional finite number>,
  "magnetic_roll": <optional finite number>,
  "beam_center": <optional local transformation>,
  "beam_length": <optional positive number>,
  "beam_curvature": <optional finite number>,
  "beam_roll": <optional finite number>,
  "frames": {
    "name": {"transformation": [...]}
  }
}
\end{lstlisting}
When present, its shape is exactly one of
\begin{align*}
  &[\texttt{"box"},d_x,d_y,d_z,\kappa,\omega],\\
  &[\texttt{"cylinder"},r,d_z,\kappa,\omega].
\end{align*}
The dimensions are positive.  The longitudinal value $d_z$ is centerline arc
length, $\kappa$ is signed constant curvature, and $\omega$ is the constant
bend-plane roll.  The type's center frame is at local path coordinate zero, so
the physical shape spans $[-d_z/2,d_z/2]$ along its curved centerline.

The four magnetic fields form one optional group: either all four are present
or all four are absent.  The four beam fields form a second independent
optional group with the same all-or-none rule.  The transformations in
\texttt{magnetic\_center}, \texttt{beam\_center}, and \texttt{frames} have no
explicit reference: their implicit reference is the type center.  An absent
group contributes no implicit frames.  An absent shape or feature group is
represented by omitted JSON members; explicit \texttt{null} values are
invalid.

Stored frame names may not use any reserved implicit frame name, whether or
not the corresponding optional group is present:
\begin{equation}
  \begin{aligned}
  &\texttt{center},\quad
   \texttt{magnetic\_center},\quad
   \texttt{magnetic\_entry},\quad
   \texttt{magnetic\_exit},\\
  &\texttt{beam\_center},\quad
   \texttt{beam\_entry},\quad
   \texttt{beam\_exit}.
  \end{aligned}
\end{equation}
The frame \texttt{center} always exists.  The three \texttt{magnetic\_*}
frames exist exactly when the magnetic group is present; the three
\texttt{beam\_*} frames exist exactly when the beam group is present.

\subsection{Objects}

An object contains only a type name and an object position:
\begin{lstlisting}
{
  "type": "typeName",
  "position": <object position>
}
\end{lstlisting}
Shape, color, magnetic and beam-interface data, and named local frames belong
to the type, not to the instance.  Any number of objects may instantiate the
same type.  A type with no shape, no magnetic group, no beam group, and no
stored frames gives each instance only its implicit \texttt{center} frame.

\section{Reference-curve geometry}
\label{sec:curve-geometry}

Consider segment $i$ with start frame
$\Frame_i=(\mathbf{r}_i,\mathbf{x}_i,\mathbf{y}_i,\mathbf{t}_i)$,
length $L_i>0$, signed angle $\alpha_i$, and roll $\omega_i$.  Its signed
curvature is
\begin{equation}
  \kappa_i=\frac{\alpha_i}{L_i}.
\end{equation}
The unit vector toward the centre of curvature and the bend axis are
\begin{align}
  \mathbf{n}_i
    &=-\cos\omega_i\,\mathbf{x}_i-
      \sin\omega_i\,\mathbf{y}_i,\\
  \mathbf{b}_i&=\mathbf{t}_i\times\mathbf{n}_i.
  \label{eq:bend-plane}
\end{align}
At zero roll, a positive segment angle bends in the local $x$--$s$ plane
toward $-\mathbf{x}$.  Increasing positive roll rotates the positive-curvature
direction from $-\mathbf{x}$ toward $-\mathbf{y}$; at $\omega_i=\pi/2$ it
points exactly toward $-\mathbf{y}$.

For a local distance $u\in[0,L_i]$ and nonzero curvature, put
$\theta=\kappa_i u$.  The exact frame is
\begin{align}
  \mathbf{r}_i(u)
    &=\mathbf{r}_i
      +\frac{\sin\theta}{\kappa_i}\mathbf{t}_i
      +\frac{1-\cos\theta}{\kappa_i}\mathbf{n}_i,
      \label{eq:arc-position}\\
  \mathbf{x}_i(u)&=\Rot(\mathbf{b}_i,\theta)\mathbf{x}_i,\\
  \mathbf{y}_i(u)&=\Rot(\mathbf{b}_i,\theta)\mathbf{y}_i,\\
  \mathbf{t}_i(u)&=\Rot(\mathbf{b}_i,\theta)\mathbf{t}_i.
  \label{eq:arc-frame}
\end{align}
Here $\Rot(\mathbf{b},\theta)$ is Rodrigues' right-handed rotation about unit
axis $\mathbf{b}$:
\begin{equation}
 \Rot(\mathbf{b},\theta)\mathbf{v}
 =\mathbf{v}\cos\theta
  +(\mathbf{b}\times\mathbf{v})\sin\theta
  +\mathbf{b}(\mathbf{b}\cdot\mathbf{v})(1-\cos\theta).
\end{equation}
For $\kappa_i=0$, the continuous limiting definition is
\begin{equation}
  \mathbf{r}_i(u)=\mathbf{r}_i+u\mathbf{t}_i,
  \qquad
  (\mathbf{x}_i(u),\mathbf{y}_i(u),\mathbf{t}_i(u))
  =(\mathbf{x}_i,\mathbf{y}_i,\mathbf{t}_i).
  \label{eq:straight-segment}
\end{equation}

Let $S_i=\sum_{j<i}L_j$ be the global path coordinate at the segment start.
The frame at $s=S_i+u$ is given by the expressions above, and the segment end
frame becomes the next segment's start frame.  This construction makes the
position and entire triad continuous at a junction.  The tangent satisfies
\begin{equation}
  \frac{d\mathbf{r}}{ds}=\mathbf{t}(s),
\end{equation}
and no independent twist is added within a segment; the triad is transported by
the common bend rotation in Equation~\eqref{eq:arc-frame}.

The defined curve domain is $[0,L_C]$, where $L_C=\sum_iL_i$.  Direct frame
evaluation is extended beyond this domain by straight tangent continuation:
\begin{align}
  \Frame_C(s)&=\Frame_C(0)T_t(s), &&s<0,\\
  \Frame_C(s)&=\Frame_C(L_C)T_t(s-L_C), &&s>L_C.
  \label{eq:curve-extrapolation}
\end{align}
Station inference, in contrast, searches only the defined domain
$[0,L_C]$.

\section{Evaluation of referenced transformations}
\label{sec:reference-transform}

Let a transformation have reference $R$ and operation list $\ops$.  Define
\begin{equation}
  \Delta s(\ops)=\sum_{(o,q)\in\ops,\;o=\texttt{ts}}q,
\end{equation}
and let $\ops_{\neg ts}$ be the same list with every \texttt{ts} removed while
preserving the relative order of all remaining operations.

For a curve reference $C$, the resolved frame is
\begin{equation}
  \operatorname{Resolve}(C,\ops)
  =\Apply\!\left(\Frame_C(\Delta s(\ops)),\ops_{\neg ts}\right).
  \label{eq:curve-reference}
\end{equation}
Thus all \texttt{ts} values on a curve reference are collected before any
other operation.  Their positions in the list do not interleave with rotations
or translations.  The non-\texttt{ts} operations still run in list order.

For a world reference, or an object-frame reference with no \texttt{ts},
\begin{align}
  \operatorname{Resolve}(\text{world},\ops)
    &=\Apply(\identity,\ops),\\
  \operatorname{Resolve}(A.f,\ops)
    &=\Apply(\Frame_{A.f}^{W},\ops).
  \label{eq:ordinary-reference}
\end{align}
An object-frame reference containing \texttt{ts} uses the special projection
semantics in Section~\ref{sec:s-inference}.

The starting frame of a curve is obtained using the same rules.  It may refer
to world, another curve, or a frame on an object.  Consequently, curves need
not be rooted directly in world coordinates.

\section{Type-local geometry, axes, and frames}
\label{sec:type-local}

\subsection{The mechanical path operator}

When type $T$ has a shape, let $\kappa_T$ and $\omega_T$ be the curvature and
roll stored in that shape.  When it has no shape, put
$\kappa_T=\omega_T=0$.  In both cases define the type-local mechanical path
operator
\begin{equation}
  \Path_T(\Frame,d)
  =\operatorname{Advance}(\Frame,d,\kappa_T d,\omega_T),
  \label{eq:type-path}
\end{equation}
where \(\operatorname{Advance}\) is the straight/arc construction of
Equations~\eqref{eq:arc-position}--\eqref{eq:straight-segment}, now permitted
for signed distance $d$.  Therefore a negative \texttt{ts} follows the same
centerline backward.  Without a shape this operator advances in a straight
line along the current tangent, so type-local \texttt{ts} remains well defined
and is geometrically identical to \texttt{tt}.

A local operation list is evaluated sequentially from the current frame:
\begin{equation}
 \Frame_{j+1}=
 \begin{cases}
   \Path_T(\Frame_j,q_j), & o_j=\texttt{ts},\\
   \Apply(\Frame_j,[(o_j,q_j)]), & o_j\ne\texttt{ts}.
 \end{cases}
 \label{eq:type-operation-recursion}
\end{equation}
Unlike a curve reference, type-local \texttt{ts} operations are \emph{not}
collected.  They execute at their exact list positions.  Hence a rotation
before \texttt{ts} changes the plane in which the local path advances, and a
\texttt{tt} before \texttt{ts} differs from a \texttt{tt} after it.

Write the resulting operator as $\Apply_T(\Frame,\ops)$.  Each elementary
operation is invertible.  The inverse list is obtained by reversing the order
and negating every value:
\begin{equation}
  \ops^{-1}=\big[(o_n,-q_n),\ldots,(o_1,-q_1)\big],
  \qquad
  \Apply_T(\Apply_T(\Frame,\ops),\ops^{-1})=\Frame.
  \label{eq:inverse-operations}
\end{equation}

\subsection{Optional swept mechanical shape}

If the type has a shape, let
\begin{equation}
  \Frame_T(u)=\Path_T(\identity,u),
  \qquad -d_z/2\le u\le d_z/2.
\end{equation}
For a box, its swept volume is parameterized by
\begin{equation}
  \mathbf{g}_{\rm box}(u,\xi,\eta)
  =\mathbf{r}_T(u)+\xi\mathbf{x}_T(u)+\eta\mathbf{y}_T(u),
  \quad
  |\xi|\le d_x/2,\quad |\eta|\le d_y/2.
  \label{eq:box-sweep}
\end{equation}
For a cylinder, a useful surface parameterization is
\begin{equation}
  \mathbf{g}_{\rm cyl}(u,\varphi)
  =\mathbf{r}_T(u)
   +r\cos\varphi\,\mathbf{x}_T(u)
   +r\sin\varphi\,\mathbf{y}_T(u),
  \quad 0\le\varphi<2\pi.
  \label{eq:cylinder-sweep}
\end{equation}
For an object with world center frame $\Frame_O^W$, the world surface follows
from Equation~\eqref{eq:local-to-world}, or equivalently by replacing
$\identity$ with $\Frame_O^W$ in the type path operator.
If the type has no shape, it contributes no surface or volume; this does not
remove its \texttt{center} frame or prevent it from being positioned.

\subsection{Stored and implicit local frames}

Every stored entry in \texttt{frames} is an operation list $\ops_f$ from the
type center.  Its local frame is
\begin{equation}
  \Frame_{T.f}=\Apply_T(\identity,\ops_f).
  \label{eq:named-local-frame}
\end{equation}
The reserved local center frame is
$\Frame_{T.\mathrm{center}}=\identity$.

For either optional feature $j\in\{m,b\}$, denoting magnetic and beam
interface respectively, define its independent path operator by
\begin{equation}
  \Path_{T,j}(\Frame,d)
  =\operatorname{Advance}(\Frame,d,\kappa_j d,\omega_j),
  \label{eq:feature-path}
\end{equation}
where $\ell_j>0$, $\kappa_j$, and $\omega_j$ are the feature's stored length,
curvature, and roll.  These values are independent of the optional shape's
mechanical path.  The finite feature axis is the locus of frame origins
$\operatorname{origin}(\Path_{T,j}(\Frame_j,u))$ for
$-\ell_j/2\le u\le\ell_j/2$.

Let $\ops_m$ be the magnetic-center transformation.  When the magnetic group
is present, its local frames are
\begin{align}
  \Frame_{T.\mathrm{magnetic\_center}}
    &=\Apply_T(\identity,\ops_m),\\
  \Frame_{T.\mathrm{magnetic\_entry}}
    &=\Path_{T,m}\!\left(
       \Frame_{T.\mathrm{magnetic\_center}},-\ell_m/2\right),\\
  \Frame_{T.\mathrm{magnetic\_exit}}
    &=\Path_{T,m}\!\left(
       \Frame_{T.\mathrm{magnetic\_center}},+\ell_m/2\right).
  \label{eq:magnetic-frames}
\end{align}

Likewise, if $\ops_b$ is the beam-center transformation and the beam group is
present,
\begin{align}
  \Frame_{T.\mathrm{beam\_center}}
    &=\Apply_T(\identity,\ops_b),\\
  \Frame_{T.\mathrm{beam\_entry}}
    &=\Path_{T,b}\!\left(
       \Frame_{T.\mathrm{beam\_center}},-\ell_b/2\right),\\
  \Frame_{T.\mathrm{beam\_exit}}
    &=\Path_{T,b}\!\left(
       \Frame_{T.\mathrm{beam\_center}},+\ell_b/2\right).
  \label{eq:beam-frames}
\end{align}

The center transformations in these equations use $\Apply_T$ and therefore
follow the mechanical type path for any \texttt{ts}, with the straight fallback
defined above when there is no shape.  In contrast, each boundary advance uses
its own feature path from its already-resolved center.  Appending a \texttt{ts}
operation to a center transformation would be incorrect whenever the feature
and mechanical curvatures or rolls differ.

Each entry or exit plane passes through the corresponding origin and is
spanned by its local $\mathbf{x}$ and $\mathbf{y}$ axes; its normal is the
local tangent $\mathbf{t}$.  Neither feature length is required to equal or
fit within the physical shape length.  The feature-center transformations may
include offsets and rotations.  If a feature group is absent, none of its
three frames or its finite axis exists.

\section{Object positioning by target alignment}
\label{sec:object-alignment}

Let object $O$ instantiate type $T$, and let its requested target frame be
$a$ (for example \texttt{center}, \texttt{magnetic\_entry}, or a stored named
frame).  First resolve the object's position specification to a desired world
target frame $\Frame_D^W$.  Then obtain the target's complete type-local frame
$\Frame_{T.a}$ using Section~\ref{sec:type-local}.  A conditional implicit
frame is a valid target only when its feature group is present.

The object center world frame is the unique frame satisfying, in homogeneous
notation,
\begin{equation}
  [\Frame_O^W][\Frame_{T.a}]=[\Frame_D^W].
  \label{eq:target-constraint}
\end{equation}
It is computed as
\begin{equation}
  \boxed{[\Frame_O^W]=[\Frame_D^W][\Frame_{T.a}]^{-1}}
  \label{eq:object-frame}
\end{equation}
This is the formal meaning of a statement such as ``place \texttt{B.start} at
a transformation from \texttt{A.end}.''  The reference and position operations
determine $\Frame_D^W$; inverting the local transform of \texttt{B.start}
determines the center of object $B$.

Once the object center is known, the world frame of any local frame $f$ is
\begin{equation}
  \boxed{[\Frame_{O.f}^W]=[\Frame_O^W][\Frame_{T.f}]}
  \label{eq:frame-world-frame}
\end{equation}
For a stored frame, $\Frame_{T.f}$ is given by
Equation~\eqref{eq:named-local-frame}; for a feature frame it is given by
Equation~\eqref{eq:magnetic-frames} or~\eqref{eq:beam-frames}.  If a shape is
present, all its points follow from Equations~\eqref{eq:box-sweep} or
\eqref{eq:cylinder-sweep} with $\Frame_O^W$ as the initial frame.

\section{Inferring a curve station from a world point}
\label{sec:s-inference}

This rule supports imported data in which a position is expressed using
\texttt{ts}, even though its primary reference is world or an object frame.
Let $P$ be the origin of that primary reference frame, and let $C$ be the named
\texttt{reference\_curve}.  The orientation of the primary reference is not
used in station inference.

\subsection{Candidate stations}

The transverse plane of the curve frame at station $s$ is
\begin{equation}
  \Pi_C(s)=
  \left\{Q\in\R^3:\big(Q-\mathbf{r}_C(s)\big)
  \cdot\mathbf{t}_C(s)=0\right\}.
\end{equation}
Candidate stations are
\begin{equation}
  \mathcal{S}(P,C)=
  \left\{s\in[0,L_C]:
  \big(P-\mathbf{r}_C(s)\big)\cdot\mathbf{t}_C(s)=0\right\}.
  \label{eq:station-equation}
\end{equation}

For a straight segment beginning at global station $S_i$, let
$\mathbf{q}=P-\mathbf{r}_i$.  The unique local candidate is
\begin{equation}
  u=\mathbf{q}\cdot\mathbf{t}_i,
  \label{eq:straight-station-root}
\end{equation}
accepted only if $0\le u\le L_i$.

For a curved segment, define $\kappa=\alpha_i/L_i$,
$\mathbf{q}=P-\mathbf{r}_i$, and
\begin{equation}
  A=\mathbf{q}\cdot\mathbf{t}_i,
  \qquad
  B=\mathbf{q}\cdot\mathbf{n}_i-\frac{1}{\kappa}.
\end{equation}
Substituting Equations~\eqref{eq:arc-position} and~\eqref{eq:arc-frame} into
Equation~\eqref{eq:station-equation} gives
\begin{equation}
  A\cos\theta+B\sin\theta=0,
  \qquad \theta=\kappa u.
  \label{eq:arc-station-root}
\end{equation}
If $(A,B)\ne(0,0)$, all roots can be enumerated as
\begin{equation}
  \theta=\operatorname{atan2}(B,A)+\frac{\pi}{2}+k\pi,
  \qquad k\in\mathbb{Z},
  \label{eq:arc-root-enumeration}
\end{equation}
retaining only those for which $u=\theta/\kappa$ is inside the segment.  If
$A=B=0$, every station in the segment is a solution; geometrically, $P$ is at
the circular arc's centre.  Identical roots at segment junctions are
de-duplicated.

\subsection{Closest-solution rule}

The set $\mathcal{S}$ may be empty, contain one or several isolated solutions,
or contain a continuous interval.  Define the distance from $P$ to a candidate
curve-frame origin by
\begin{equation}
  d(s)=\lVert P-\mathbf{r}_C(s)\rVert_2.
\end{equation}
The inferred station $\widehat{s}$ exists if and only if there is one unique,
isolated minimizer of $d$ over $\mathcal{S}$:
\begin{equation}
  \widehat{s}
  =\underset{s\in\mathcal{S}(P,C)}{\operatorname{argmin}}\;d(s).
  \label{eq:closest-station}
\end{equation}
Accordingly:
\begin{itemize}
  \item if $\mathcal{S}$ is empty, evaluation raises an exception;
  \item if one isolated solution is strictly closest, it is selected even when
        farther solutions or a farther continuous interval exist;
  \item if two or more closest solutions are equidistant, evaluation raises an
        ambiguity exception;
  \item if a closest continuous interval exists, evaluation raises the same
        ambiguity exception.
\end{itemize}
Floating-point implementations shall use scale-aware distance and path
tolerances both for equality and for junction de-duplication.

\subsection{Applying the inferred station}

Let $\Delta s$ be the sum of the object's \texttt{ts} values and
$\ops_{\neg ts}$ its remaining operations.  The desired target frame is
\begin{equation}
  \boxed{
  \Frame_D^W
  =\Apply\!\left(\Frame_C(\widehat{s}+\Delta s),
                  \ops_{\neg ts}\right)}.
  \label{eq:inferred-position}
\end{equation}
The origin $P$ is used only to infer $\widehat{s}$.  Both the origin and
orientation of the primary reference frame are then replaced by the curve frame
at $\widehat{s}+\Delta s$.  The remaining operations are applied in their
original relative order.  Finally Equation~\eqref{eq:object-frame} aligns the
object's requested target frame with $\Frame_D^W$.

\section{Complete world-resolution algorithm}
\label{sec:resolution-algorithm}

World geometry is obtained by recursively resolving named curves and objects.
The following algorithm is normative at the semantic level.

\begin{enumerate}
  \item Validate the canonical structure, finite numbers, positive dimensions
        and lengths, colors, operation names, target names, and all references.
  \item Construct a dependency graph whose nodes are named curves and objects.
        A curve depends on the curve or object named by its starting-frame
        reference.  An object depends on the curve or object named by its
        position reference; a non-curve object position containing \texttt{ts}
        additionally depends on its \texttt{reference\_curve}.
  \item Reject the layout if that graph contains a directed cycle.
  \item Resolve nodes in topological order, or use memoized recursion with an
        active-resolution stack.
  \item To resolve a curve, resolve its starting frame using
        Equations~\eqref{eq:curve-reference} or
        \eqref{eq:ordinary-reference}, then propagate all segments using
        Section~\ref{sec:curve-geometry}.
  \item To resolve an object, resolve its desired target frame:
        \begin{itemize}
          \item a curve reference uses Equation~\eqref{eq:curve-reference};
          \item a world or object-frame reference without \texttt{ts} uses
                Equation~\eqref{eq:ordinary-reference};
          \item a world or object-frame reference with \texttt{ts} uses
                Equations~\eqref{eq:closest-station} and
                \eqref{eq:inferred-position}.
        \end{itemize}
  \item Align the object's target using Equation~\eqref{eq:object-frame}.
  \item Resolve its stored, magnetic, and beam-interface frames that actually
        exist using Equation~\eqref{eq:frame-world-frame}.  If the type has a
        shape, sweep it in world coordinates using
        Section~\ref{sec:type-local}.
\end{enumerate}

The following identities summarize the output of the algorithm:
\begin{align}
  \text{object centre:}\quad
  [\Frame_O^W]&=[\Frame_D^W][\Frame_{T.a}]^{-1},\\
  \text{named frame:}\quad
  [\Frame_{O.f}^W]&=[\Frame_O^W][\Frame_{T.f}],\\
  \text{world shape point:}\quad
  \begin{bmatrix}\mathbf{p}_W\\1\end{bmatrix}
  &=[\Frame_O^W]
    \begin{bmatrix}\mathbf{p}_T\\1\end{bmatrix}.
\end{align}

\section{Validity and failure conditions}

A conforming reader rejects a layout when any of the following holds:
\begin{itemize}
  \item a JSON object contains unsupported fields or has a missing required
        field;
  \item an entity name is empty, an operation value is not finite, a curve
        segment length is not positive, a present shape dimension is not
        positive, or a present magnetic or beam length is not positive;
  \item a present shape is not exactly one of the two canonical forms
        \begin{align*}
          &[\texttt{"box"},d_x,d_y,d_z,\kappa,\omega],\\
          &[\texttt{"cylinder"},r,d_z,\kappa,\omega];
        \end{align*}
  \item a color is not a six-digit hexadecimal color;
  \item only some fields of the magnetic group or only some fields of the beam
        group are present;
  \item a present shape, magnetic, or beam curvature or roll is not finite;
  \item a type, curve, object, stored frame, applicable implicit frame, or
        positioning target does not exist;
  \item a stored frame uses any reserved implicit name, including a name whose
        optional feature is absent;
  \item \texttt{ts} is used with a world or object-frame reference outside an
        object position, or it is used in an object position without the
        required \texttt{reference\_curve};
  \item \texttt{reference\_curve} is supplied with a primary curve reference;
  \item the reference dependency graph contains a cycle;
  \item station inference has no solution or has equidistant closest solutions,
        including a closest continuous interval.
\end{itemize}

\section{World-pose readout and MAD-X-style angles}
\label{sec:madx-angles}

The world position of any resolved frame is simply its origin
$\mathbf{o}=(X,Y,Z)$.  For display, the current editor derives
MAD-X-style survey angles from the frame as follows.  Let
\begin{align}
 \widehat{\mathbf{t}}&=\frac{\mathbf{t}}{\lVert\mathbf{t}\rVert},\\
 \widehat{\mathbf{x}}&=
 \frac{\mathbf{x}-(\mathbf{x}\cdot\widehat{\mathbf{t}})
       \widehat{\mathbf{t}}}
      {\left\lVert\mathbf{x}-(\mathbf{x}\cdot\widehat{\mathbf{t}})
       \widehat{\mathbf{t}}\right\rVert},\\
 \widehat{\mathbf{y}}&=\widehat{\mathbf{t}}\times
                       \widehat{\mathbf{x}},\\
 h&=\sqrt{t_X^2+t_Z^2}.
\end{align}
Away from the vertical singularity,
\begin{align}
 \theta&=\operatorname{atan2}(t_X,t_Z),\\
 \phi&=\operatorname{atan2}(t_Y,h),\\
 \psi&=\operatorname{atan2}(x_Y,y_Y).
 \label{eq:madx-readout}
\end{align}
Here $\theta$ is horizontal azimuth, $\phi$ is vertical elevation, and $\psi$
is roll about the local tangent.  At $h\approx0$, the display fixes
$\theta=0$ and uses
$\psi=\operatorname{atan2}(-y_X,x_X)$.  This branch choice makes the readout
deterministic at the usual Euler-angle singularity; the frame matrix remains the
authoritative pose.

\section{Canonical example}

The following layout places \texttt{Q1.center} at $s=10$ m on a straight
reference curve.  The reusable type has a mechanical shape, independent
magnetic and beam-interface axes, and two stored local frames.  All three axes
are straight in this compact example.

\begin{lstlisting}
{
  "reference_curves": {
    "main": {
      "color": "#7d91ff",
      "starting_frame": {
        "reference": {"kind": "world"},
        "transformation": []
      },
      "segments": [[100.0, 0.0, 0.0]]
    }
  },
  "types": {
    "magnet": {
      "shape": ["box", 2.0, 2.0, 2.0, 0.0, 0.0],
      "color": "#f0a84b",
      "magnetic_center": {"transformation": []},
      "magnetic_length": 2.0,
      "magnetic_curvature": 0.0,
      "magnetic_roll": 0.0,
      "beam_center": {"transformation": []},
      "beam_length": 2.0,
      "beam_curvature": 0.0,
      "beam_roll": 0.0,
      "frames": {
        "start": {"transformation": [["ts", -1.0]]},
        "end":   {"transformation": [["ts",  1.0]]}
      }
    }
  },
  "objects": {
    "Q1": {
      "type": "magnet",
      "position": {
        "target": "center",
        "reference": {"kind": "curve", "curve": "main"},
        "transformation": [["ts", 10.0]]
      }
    }
  }
}
\end{lstlisting}

For this example,
\begin{align*}
 \Frame_{Q1.\mathrm{center}}^W&=\Frame_{\mathrm{main}}(10),\\
 \Frame_{Q1.\mathrm{start}}^W
   &=\Path_{\mathrm{magnet}}(\Frame_{Q1.\mathrm{center}}^W,-1),\\
 \Frame_{Q1.\mathrm{end}}^W
   &=\Path_{\mathrm{magnet}}(\Frame_{Q1.\mathrm{center}}^W,+1).
\end{align*}
The magnetic and beam entry frames also have origin $(0,0,9)$ m, and their
exit frames have origin $(0,0,11)$ m.  Since every curvature is zero, the
center, stored-frame, magnetic, and beam-frame orientations are identical.
In particular, the three stored-frame origins above are
$(0,0,10)$ m, $(0,0,9)$ m, and $(0,0,11)$ m respectively, with identical
orientations.

\section{Conformance notes}

The homogeneous frame is the authoritative result.  Sampled polylines,
triangulated display meshes, transparent rendering, curve tubes, hover frames,
and station snapping are viewer concerns and do not change the equations in
this specification.  In particular, implementations should evaluate an arc
analytically for positioning even if they draw it using samples.

Named frames carry both an origin and an orientation because alignment normally
requires both.  Magnetic and beam-interface entry and exit locations are
likewise frames; their $x$--$y$ planes are the physically meaningful
transverse planes.  A viewer may independently hide the magnetic axis and its
entry/exit frames, the beam-interface axis and its entry/exit frames, and
stored named frames; these three layers are off by default.
Visibility does not change
which frames exist or how they resolve.  A later beam-physics export may use
the beam-interface frames to construct element boundaries and derive survey
offsets or misalignments without changing the installation layout itself.

\end{document}
